\chapter{TRABALHOS RELACIONADOS}\label{chap_trabalhos_relacionados}

Esta seção tem como objetivo analisar o estado da arte no campo do aprendizado de máquina quântico, enfocando tanto a teoria quanto as aplicações práticas. As referências discutidas aqui foram selecionadas com base em sua relevância e contribuição para o desenvolvimento do tema proposto.

\section{Computação Quântica}

\citeonline{Mermin2007QComputerScience} oferece uma introdução abrangente à computação quântica, discutindo conceitos como qubits e cbits, preparação e medição de estados quânticos, e a solução de problemas como Deutsch, Bernstein-Vazirani e Simon. O livro também explora a criptografia RSA e como quebrá-la, busca em computadores quânticos e algoritmos como a iteração de Grover, além de abordar a correção de erros quânticos.

\citeonline{Wong2022IntroductionClassicalQuantum} apresenta uma introdução simplificada à computação clássica e quântica, abordando conceitos básicos da computação clássica, incluindo bits, portas lógicas e programação de circuitos clássicos usando linguagens de descrição de hardware. Ele também discute as classes de problemas que são fáceis ou difíceis para os computadores e a perspectiva de que os computadores quânticos possam ser significativamente mais rápidos do que os computadores clássicos em algumas tarefas. O artigo então aborda os conceitos básicos da computação quântica, incluindo qubits, superposições e portas quânticas. O artigo enfatiza a importância da álgebra linear na computação quântica e mostra como ela pode ser usada. Em última análise, o artigo tem como objetivo fornecer uma introdução ao campo da computação quântica para iniciantes.

\citeonline{Zubairy2020} visa apresentar os conceitos básicos da mecânica quântica para iniciantes e aplicá-los a várias aplicações, como comunicação quântica e computação quântica. O artigo também discute os fundamentos da mecânica quântica e sua relação com a dinâmica newtoniana. O artigo tem como objetivo fornecer uma compreensão abrangente da mecânica quântica e suas aplicações.

\section{Aprendizado de Máquina Quântico}

\citeonline{Biamonte2017QML} explora a interseção entre computação quântica e aprendizado de máquina e como eles podem ser usados juntos para aprimorar algoritmos de aprendizado e fornecer \textit{speedups}. O artigo discute vários exemplos de como a mecânica quântica pode ser usada para melhorar os métodos de aprendizagem clássicos e como os algoritmos quânticos podem ser usados para várias tarefas. Os autores também discutem o potencial da computação quântica e do aprendizado de máquina coevoluírem e se tornarem tecnologias facilitadoras uma para a outra.

\citeonline{Khan2020ML} fornece uma visão geral do aprendizado de máquina quântico em comparação com as abordagens clássicas. O artigo discute as contribuições técnicas, pontos fortes e semelhanças do trabalho de pesquisa neste domínio. Ele também aborda o progresso recente em diferentes abordagens de aprendizado de máquina quântica, sua complexidade e aplicações em vários campos, como física, química e processamento de linguagem natural. O artigo tem como objetivo explorar o uso de sistemas quânticos para processar dados clássicos usando algoritmos de aprendizado de máquina e como algoritmos e softwares de aprendizado de máquina podem ser formulados e implementados em computadores quânticos. O artigo também destaca os desafios significativos de hardware e software que precisam ser resolvidos antes que o aprendizado de máquina quântico se torne prático.

\citeonline{Havlicek2019qSupervisedLearning} apresenta um novo método para aprendizado supervisionado usando espaços de recursos quânticos aprimorados. Os autores propõem uma abordagem baseada em kernel quântico que pode ser usada para classificar conjuntos de dados com alta precisão. Eles demonstram a eficácia de seu método aplicando-o a vários conjuntos de dados de referência e comparando os resultados com algoritmos clássicos de aprendizado de máquina.

\citeonline{Rebentrost2014QSVM} demonstra que a máquina vetorial de suporte, um método de classificação popular em aprendizado de máquina, pode ser implementada em um computador quântico com complexidade logarítmica no tamanho do recurso e no número de dados de treinamento. O artigo mostra que esse algoritmo quântico pode fornecer um \textit{speedup} exponencial nos casos em que algoritmos de amostragem clássicos exigem tempo polinomial. O artigo também destaca os benefícios do aprendizado de máquina quântica em termos de privacidade de dados e sugere que a máquina vetorial de suporte quântico poderia ser usada como um componente em uma rede neural quântica maior.

\citeonline{Zhou2018HandwrittenDigitRecognition} estabelece um modelo de aprendizado de máquina de reconhecimento de dígitos manuscritos baseado no \textit{Support Vector Machine} e analisar a influência do número da amostra, dos parâmetros da função do kernel, dos coeficientes de penalidade e de outros parâmetros no modelo de predição. O estudo visa melhorar o desempenho do aprendizado de máquina no reconhecimento de informações numéricas manuscritas.

O livro de \citeonline{Schuld2021MLQComputers} fornece uma introdução e um contexto para o aprendizado de máquina quântico. O objetivo é orientar os leitores sobre diferentes maneiras de pensar sobre o aprendizado de máquina a partir de uma perspectiva de computação quântica. O livro também apresenta um exemplo de como os computadores quânticos podem aprender com os dados, que exibe uma série de questões discutidas no decorrer do livro.

\citeonline{Métricas} discute várias métricas de avaliação que os cientistas de dados podem usar para avaliar seus modelos de classificação e fornecer orientação sobre como e quando usá-los. O artigo enfatiza que diferentes problemas de negócios exigem estratégias de otimização diferentes e que a precisão nem sempre é a melhor métrica a ser usada para avaliar o desempenho do modelo.

\section{Conclusão do Capítulo}

Os trabalhos relacionados descritos nesta seção foram usados para adquirir conhecimento teórico sobre física quântica, computação quântica e métodos de aprendizado de máquina quântico, bem como suas aplicações. A maioria das referências aborda o tema do \textit{quantum machine learning}, mas também inclui dois livros que oferecem uma introdução à computação quântica e esclarecem suas bases físicas.